\section{A First Lean 4 Reading}
\label{sec:lean-4-first-reading}

\begin{topicbox}{Subsets as Predicates}
\begin{exposition}
Lean represents many mathematical objects by their defining behavior. For a
type $P$, a subset of $P$ can be read as a predicate $P\to Prop$: it takes an
element of $P$ as input and returns the proposition that the element belongs to
the subset. This is the key syntax bridge for the Peano-system arguments that
use inductive subsets.
\end{exposition}

\begin{figure}[H]
\centering
\resizebox{\textwidth}{!}{%
\begin{tikzpicture}[
  node distance=1.2cm,
  code/.style={
    draw,
    rounded corners,
    fill=gray!8,
    font=\ttfamily\small,
    align=left,
    inner sep=8pt
  },
  note/.style={
    draw,
    rounded corners,
    fill=blue!6,
    font=\small,
    align=left,
    text width=4.2cm,
    inner sep=6pt
  },
  arrow/.style={->, thick}
]

\node[code] (code) {
let D : ps.carrier $\to$ Prop :=\\
\quad fun candidate\_element =>\\
\quad\quad candidate\_element = ps.one $\lor$\\
\quad\quad $\exists$ predecessor : ps.carrier,\\
\quad\quad\quad ps.successor predecessor = candidate\_element
};

\node[note, right=2.8cm of code.north east, anchor=north west] (subset) {
\textbf{D is a subset predicate}\\
In Lean, a subset of the carrier is represented as a function:
\[
D : ps.carrier \to Prop.
\]
It answers: ``does this element belong to \(D\)?''
};

\node[note, right=2.8cm of code.east, anchor=west] (candidate) {
\textbf{candidate\_element}\\
This is the element currently being tested for membership in \(D\).
It plays the role of \(x\) in:
\[
D=\{x\in P : \cdots\}.
\]
};

\node[note, right=2.8cm of code.south east, anchor=south west] (condition) {
\textbf{Membership condition}\\
The proposition says:
\[
x = 1
\quad\text{or}\quad
\exists u,\ S(u)=x.
\]
So \(D\) contains exactly elements that are either \(1\) or successors.
};

\node[note, below=1.5cm of code] (predecessor) {
\textbf{predecessor}\\
The existential witness:
\[
\exists\, predecessor : ps.carrier
\]
means ``there is some element of the carrier whose successor is the candidate element.''
};

\draw[arrow] (subset.west) -- ++(-0.9,0) |- ([yshift=0.45cm]code.east);
\draw[arrow] (candidate.west) -- ++(-0.9,0) |- ([yshift=-0.15cm]code.east);
\draw[arrow] (condition.west) -- ++(-0.9,0) |- ([yshift=-0.95cm]code.east);
\draw[arrow] (predecessor.north) -- ([yshift=-1.2cm]code.south);

\end{tikzpicture}

}
\caption{Reading a Lean subset predicate as a mathematical subset condition.}
\label{fig:lean-function-syntax}
\end{figure}

\begin{remark*}[Interpretation]
The Lean line `D : ps.carrier \(\to\) Prop` says that membership in $D$ is a
proposition depending on an input from the carrier of the Peano system. In
ordinary mathematical notation, this is the same role played by a subset
$D\subseteq P$ or by a defining condition $D(x)$.
\end{remark*}
\end{topicbox}

\begin{topicbox}{Reading an Initial Proof}
\begin{exposition}
Lean proofs are built from local goals and the terms that solve them. A first
reading strategy is to identify the current hypotheses, the target proposition,
and the exact line that transforms the local context into the required result.
\end{exposition}

\begin{figure}[H]
\centering
\resizebox{\textwidth}{!}{%
\begin{tikzpicture}[
  code/.style={
    draw,
    rounded corners,
    fill=gray!8,
    font=\ttfamily\small,
    align=left,
    inner sep=8pt
  },
  note/.style={
    draw,
    rounded corners,
    fill=blue!6,
    font=\small,
    align=left,
    text width=4.4cm,
    inner sep=6pt
  },
  arrow/.style={->, thick}
]

\node[code] (theorem) {
theorem every\_element\_is\_one\_or\_successor\\
\quad (ps : PeanoSystem) :\\
\quad $\forall$ element : ps.carrier,\\
\quad\quad element = ps.one $\lor$\\
\quad\quad $\exists$ predecessor : ps.carrier,\\
\quad\quad\quad ps.successor predecessor = element := by\\[2mm]
\quad apply induction\_principle ps\\[1mm]
\quad $\cdot$ left\\
\quad\quad rfl\\[1mm]
\quad $\cdot$ intro element induction\_hypothesis\\
\quad\quad right\\
\quad\quad exact $\langle$element, rfl$\rangle$
};

\node[note, right=2.4cm of theorem.north east, anchor=north west] (ps) {
\textbf{\texttt{ps : PeanoSystem}}\\
The theorem is not about one fixed natural number system. It is about an arbitrary Peano system:
\[
(P,S,1).
\]
};

\node[note, right=2.4cm of theorem.east, anchor=west] (goal) {
\textbf{The target predicate}\\
Lean infers the predicate:
\[
P(x) := x=1 \;\lor\; \exists u,\; S(u)=x.
\]
This is the Lean-native version of defining the subset \(D\).
};

\node[note, right=2.4cm of theorem.south east, anchor=south west] (base) {
\textbf{Base case}\\
For \(x=1\), prove:
\[
1=1 \;\lor\; \exists u,\; S(u)=1.
\]
Choose the left side using \texttt{left}; prove \(1=1\) by \texttt{rfl}.
};

\node[note, below=1.3cm of theorem] (step) {
\textbf{Successor step}\\
To prove the statement for \(S(x)\), choose the right side:
\[
\exists u,\; S(u)=S(x).
\]
Take \(u=x\). Then \(S(x)=S(x)\), proved by \texttt{rfl}.
};

\node[note, left=2.4cm of theorem.west, anchor=east] (ih) {
\textbf{Induction hypothesis}\\
The variable \texttt{induction\_hypothesis} says the theorem already holds for \texttt{element}.\\
In this proof, it is not actually needed, because \(S(\texttt{element})\) is visibly a successor.
};

\draw[arrow] (ps.west) -- ++(-0.7,0) |- ([yshift=2.1cm]theorem.east);
\draw[arrow] (goal.west) -- ++(-0.7,0) |- ([yshift=0.6cm]theorem.east);
\draw[arrow] (base.west) -- ++(-0.7,0) |- ([yshift=-1.0cm]theorem.east);
\draw[arrow] (step.north) -- ([yshift=-1.7cm]theorem.south);
\draw[arrow] (ih.east) -- ++(0.7,0) |- ([yshift=-1.25cm]theorem.west);

\end{tikzpicture}
}
\caption{Reading the shape of an initial Lean proof.}
\label{fig:lean-initial-proof}
\end{figure}

\begin{remark*}[Interpretation]
The figure isolates the proof-state reading habit used throughout Lean:
separate the available assumptions from the current goal, then read each proof
term as an instruction that closes or transforms that goal.
\end{remark*}
\end{topicbox}
